% 2.6 — brief en documento/investigacion/2-6-ingenieria-de-software.md
\subsection{INGENIERÍA DE SOFTWARE}
\subsubsection{PROCESO DE DESARROLLO DE SOFTWARE}

Un proceso de desarrollo de software es el conjunto de actividades que conduce a la producción de un producto de software. La norma \textcite{ISO12207-2017} establece un marco común de procesos del ciclo de vida, con terminología definida, que contiene procesos, actividades y tareas aplicables durante la adquisición, el suministro, el desarrollo, la operación, el mantenimiento y el retiro de sistemas, productos y servicios de software; los agrupa en procesos de acuerdo, de habilitación organizacional, de gestión técnica y técnicos. La norma aclara, sin embargo, que no prescribe ningún modelo de ciclo de vida ni metodología de desarrollo: quien la usa es responsable de elegir un modelo para su proyecto y de mapear sobre él los procesos y actividades. Esa distinción entre proceso (qué actividades se realizan) y modelo de ciclo de vida (cómo se ordenan en el tiempo) es la que organiza esta sección y la siguiente.

\textcite{Sommerville2016} presenta los modelos de proceso como representaciones abstractas y distingue tres: el modelo en cascada, dirigido por plan, con fases separadas de especificación y desarrollo; el desarrollo incremental, que intercala especificación, desarrollo y validación; y la integración y configuración, basada en reutilizar componentes existentes. \textcite{Pressman2020} agrupan los primeros bajo la etiqueta de modelos prescriptivos, en contraste con los procesos ágiles. El modelo en cascada suele remontarse a \textcite{Royce1970}, que propone una secuencia de requisitos de sistema, requisitos de software, análisis, diseño de programa, codificación, pruebas y operación, aunque el nombre es posterior y el propio autor advierte que esa implementación es riesgosa e invita al fracaso, porque la fase de pruebas al final es el primer momento en que se experimentan, y no solo se analizan, los tiempos y el almacenamiento, y un fallo ahí obliga a rediseñar. Sus correcciones incluyen documentar el diseño, hacerlo dos veces (que la versión entregada sea en realidad la segunda) e involucrar al cliente.

Esa crítica se formaliza en el modelo en espiral. \textcite{Boehm1988} lo presenta como un enfoque evolutivo guiado por el riesgo que provee un marco para conducir el proceso de software; según \textcite{Larman2003}, su aporte fue hacer prominente el concepto de iteraciones guiadas por riesgo, con un paso explícito de evaluación de riesgos en cada iteración. Los métodos ágiles llevan la iteración más lejos: el manifiesto de \textcite{Beck2001} antepone los individuos y sus interacciones a los procesos y herramientas, el software funcionando a la documentación exhaustiva, la colaboración con el cliente a la negociación contractual y la respuesta al cambio al seguimiento de un plan. \textcite{Larman2003} señalan que los diecisiete firmantes representaban métodos como XP, Scrum y DSDM, todos ellos iterativos e incrementales. El cuadro~\ref{tab:modelos-proceso} resume los modelos comparados.

\begin{table}[H]
  \centering
  \caption{Modelos de proceso de software comparados}\label{tab:modelos-proceso}
  \begin{tabular}{@{}>{\RaggedRight\arraybackslash}p{0.14\textwidth}>{\RaggedRight\arraybackslash}p{0.30\textwidth}>{\RaggedRight\arraybackslash}p{0.24\textwidth}>{\RaggedRight\arraybackslash}p{0.24\textwidth}@{}}
    \toprule
    \textbf{Modelo} & \textbf{Rasgo central} & \textbf{Fortaleza} & \textbf{Limitación} \\
    \midrule
    Cascada & Fases secuenciales con especificación cerrada antes del desarrollo & Simple de explicar; proceso visible y medible & Los problemas se descubren en las pruebas finales; el cambio es costoso \\
    Espiral & Ciclos guiados por el riesgo, con evaluación explícita en cada uno & Reduce el riesgo antes de comprometer recursos & Exige experiencia en análisis de riesgos \\
    Incremental & Especificación, desarrollo y validación intercalados en entregas parciales & Menor costo del cambio; retroalimentación temprana; entrega útil más rápida & Menor visibilidad para la gestión; la estructura se degrada sin refactorización \\
    Ágil & Iteraciones cortas, colaboración continua con el cliente, respuesta al cambio & Adaptación rápida; software funcionando como medida de avance & Depende de la disponibilidad del cliente y de equipos autoorganizados \\
    \bottomrule
  \end{tabular}
  \fuente{Elaboración propia a partir de \textcite{Royce1970}, \textcite{Boehm1988}, \textcite{Sommerville2016}, \textcite{Larman2003} y \textcite{Beck2001}.}
\end{table}

Para este trabajo, las actividades del proceso quedan fijas: cada incremento comprende definición de alcance, análisis, diseño, desarrollo y validación, que corresponden a los procesos técnicos de la norma. Lo que se elige es el modelo de ciclo de vida que las ordena, y se prevé que sea el iterativo e incremental, por las razones que se desarrollan en la subsección siguiente.

\subsubsection{METODOLOGÍA DE DESARROLLO ITERATIVA E INCREMENTAL}

\textcite{Larman2003} reconstruyen la historia del desarrollo iterativo e incremental y muestran que, aunque muchos lo consideran una práctica moderna, su aplicación se remonta a mediados de los años cincuenta, con un ejemplo temprano en el Proyecto Mercury de la NASA, que trabajó con iteraciones de medio día. La definición clásica proviene de \textcite{Basili1975}, citados por \textcite{Larman2003}: comenzar con una implementación simple de un subconjunto de los requisitos y mejorar iterativamente la secuencia de versiones hasta implementar el sistema completo, incorporando en cada iteración modificaciones de diseño junto con nuevas capacidades funcionales. El rasgo común de todos los enfoques de esta familia es evitar el paso único, secuencial, dirigido por documentos y con compuertas entre fases; y la palabra «iterativo» implica no solo revisar el trabajo, sino avanzar evolutivamente.

La misma fuente documenta proyectos que fijaron de antemano el número de incrementos: el sistema de defensa de TRW para el ejército estadounidense se desarrolló en cinco iteraciones en 1972; y el proyecto Trident de IBM, en cuatro iteraciones de unos seis meses. En el plano normativo, la norma militar Mil-Std-498 de 1994 reemplazó la exigencia de cascada de una sola pasada por un desarrollo en una o más construcciones incrementales, cada una de las cuales implementa un subconjunto especificado de las capacidades planificadas. \textcite{Sommerville2016} resume los beneficios del desarrollo incremental frente a la cascada: menor costo de acomodar cambios en los requisitos, retroalimentación más fácil del cliente y entrega más rápida de software útil; y sus problemas: menor visibilidad para la gestión y degradación de la estructura si no se refactoriza.

Tres condiciones del proyecto justifican la elección del modelo iterativo e incremental sobre la cascada. La primera es que el desarrollo lo realiza una sola persona: un modelo secuencial concentraría toda la validación al final, que es exactamente el riesgo que \textcite{Royce1970} describe. La segunda es que el alcance se fija por incrementos: cada uno implementa un subconjunto especificado de las capacidades planificadas, en la lógica de las construcciones incrementales de Mil-Std-498, y cierra con una validación propia, de modo que al terminar cada uno existe software funcionando. La tercera es la dependencia de los datos: las variables predictoras y los casos de uso solo se confirman al explorar el histórico del portal anterior, lo que impide cerrar la especificación completa de antemano; el modelo incremental es el que abarata ese cambio. Como la norma \textcite{ISO12207-2017} deja la elección del modelo de ciclo de vida al proyecto, esta decisión no contradice el marco de procesos.

Se prevé, en consecuencia, un desarrollo en cinco incrementos: los procesos del portal anterior migrados al ERP, la autenticación centralizada con Keycloak, el pipeline de datos y el entrenamiento de los modelos, el microservicio de predicción con su integración al ERP y al tablero, y la integración con Power BI y MCP\@. La documentación de cada incremento en el Capítulo IV mitiga el problema de visibilidad que señala Sommerville; la degradación de la estructura se atiende con la arquitectura por capas que se describe a continuación.

\subsubsection{ARQUITECTURA DE SOFTWARE POR CAPAS}

\textcite{Bass2021} entienden la arquitectura de software como el conjunto de estructuras necesarias para razonar sobre el sistema, compuestas por elementos, relaciones entre ellos y propiedades de ambos, y distinguen las estructuras de módulos, entre ellas las capas, de las de componentes y conectores y de las de asignación al entorno de despliegue. \textcite{Richards2020} añaden que las decisiones de arquitectura definen las reglas de cómo debe construirse el sistema: por ejemplo, que solo las capas de negocio y de servicios puedan acceder a la base de datos, con lo que se impide que la capa de presentación haga llamadas directas.

El patrón en capas es una estructura de módulos. \textcite{Sommerville2016} lo describe como la organización del sistema en capas, cada una con funcionalidad relacionada, donde cada capa provee servicios a la capa superior y las inferiores representan los servicios centrales; su ventaja principal es que permite reemplazar una capa entera mientras se mantenga su interfaz, y sus desventajas son la dificultad de lograr una separación limpia y el costo de rendimiento de atravesar varios niveles para atender una petición. \textcite{Richards2020} formulan el principio de capas de aislamiento, según el cual una capa depende solo de la inmediata inferior. \textcite{Fowler2002} fija el vocabulario para aplicaciones empresariales con tres capas principales, presentación, lógica de dominio y fuente de datos, y señala que esas capas lógicas pueden ejecutarse en distintos nodos físicos.

Cuando la separación se lleva más allá de un proceso, se llega a los microservicios. \textcite{Lewis2014} los definen como un enfoque para desarrollar una única aplicación como un conjunto de servicios pequeños, cada uno ejecutándose en su propio proceso y comunicándose mediante mecanismos ligeros, a menudo una API HTTP de recursos, y entre sus características cuentan la organización alrededor de capacidades de negocio, la gestión descentralizada de datos y la gobernanza descentralizada, que admite distintas tecnologías por servicio. \textcite{Newman2021} trata la extracción de capacidades de un monolito hacia servicios, la comunicación síncrona por HTTP entre ellos y la autenticación entre servicios. \textcite{Richards2020} ponen capas y microservicios en el mismo marco comparativo: el estilo distribuido aporta contexto acotado, aislamiento de datos y una capa de API, a cambio de costos operativos.

En este trabajo se combinan ambos niveles. Dentro de cada servicio se aplica el patrón en capas: en el ERP, controladores, servicios de dominio y acceso a datos con Prisma sobre PostgreSQL, y Angular como presentación; en el módulo predictivo, rutas, lógica de predicción y acceso a modelos y datos. Entre servicios, el módulo predictivo se construye como un microservicio en FastAPI separado del ERP en NestJS: es una decisión de diseño justificada porque corre en su propio proceso y expone una API HTTP, es dueño de sus modelos y de sus datos analíticos, usa un lenguaje distinto al del ERP en la lógica de la gobernanza descentralizada, y puede reentrenarse y desplegarse sin tocar el ERP, lo que responde a la modificabilidad y la integrabilidad que \textcite{Bass2021} tratan como atributos de calidad. Todos los servicios validan la identidad contra Keycloak y se despliegan en contenedores Docker en un servidor propio de la empresa. Se prevé que esta organización se concrete en la arquitectura de seis capas del módulo predictivo, cuyo diseño detallado corresponde a la sección~3.4.
